\section{Online Learning: una posible ruta de nivel paradigmático}

Las secciones anteriores explicaron por qué un Agent necesita un entorno y retroalimentación; esta sección aborda el Online Learning. Guangmi (广密) considera que, si en el futuro surge otra ruta nueva de nivel paradigmático, el Online Learning es una de las candidatas. Su núcleo no es entrenar el modelo y después desplegarlo, sino que el modelo explore de forma autónoma en un entorno, obtenga retroalimentación, actualice su experiencia e incluso, más adelante, actualice sus pesos o su memoria de largo plazo.

\begin{figure}[H]
\centering
\includegraphics[width=0.96\textwidth]{online-learning-loop.png}
\caption{Ciclo cerrado del Online Learning: el modelo explora de forma autónoma en el entorno, obtiene retroalimentación y actualiza sus capacidades. Diagrama conceptual de elaboración propia, basado en la conversación de 00:55:49--01:02:45.}
\end{figure}

\begin{knowledgebox}{Lectura de la figura: el Online Learning incorpora la producción de datos al proceso de uso}
De Environment a Action, Reward, Update y Memory/Weights, la figura subraya un ciclo continuo. El entrenamiento tradicional se parece más al aprendizaje fuera de línea; el Online Learning, en cambio, hace que el modelo aprenda mientras actúa en un entorno real o simulado. Cuantos más Agents haya, más datos de retroalimentación potenciales habrá.
\end{knowledgebox}

\subsection{Por qué el Online Learning podría ser el próximo paradigma}

Esta sección explica su atractivo. El Pre-training depende principalmente de corpus ya existentes, y el Post-training/RL depende de tareas y recompensas ya organizadas; el Online Learning, en cambio, podría generar datos nuevos a partir de las propias acciones del modelo. En entornos como Coding, videojuegos, experimentos científicos u operación de páginas web, el modelo puede intentar, fallar, corregir y sintetizar, y luego reincorporar esa experiencia. Esto desplaza el cuello de botella de los datos de «encontrar más datos antiguos» a «crear experiencia nueva de mayor valor».

\begin{importantbox}{Cambio central: del conjunto de datos al entorno}
Si el Online Learning se confirma, el centro del entrenamiento pasará de los conjuntos de datos estáticos a los entornos interactivos. Quien posea entornos, tareas, recompensas, evaluaciones y retroalimentación de usuarios reales podría tener un nuevo volante de datos.
\end{importantbox}

\subsection{Impacto en la GPU y en la narrativa de NVIDIA}

Esta sección responde a la pregunta sobre la industria que se planteó en la entrevista. El Online Learning no necesariamente debilita la demanda de GPU; por el contrario, podría cambiar su estructura: no solo se necesita entrenar, sino también muestrear, ejecutar, hacer inferencia de múltiples turnos, evaluar y reproducir a gran escala. La inferencia y el entrenamiento podrían entrelazarse más estrechamente, y la infra tendría que soportar al mismo tiempo rollout en línea, caché, simulación de entornos y auditoría de seguridad. La narrativa de NVIDIA tampoco se limita a vender tarjetas para entrenamiento, sino que se amplía en torno a la inferencia, la simulación, la robótica, los centros de datos y la pila de software.

\begin{knowledgebox}{Términos clave: infraestructura relacionada con el Online Learning}
\noindent\begin{tabularx}{\textwidth}{p{0.20\textwidth}p{0.34\textwidth}X}
\toprule
Término & Significado & Por qué importa \\
\midrule
Rollout & El Agent ejecuta una trayectoria en el entorno & Produce experiencia nueva que puede evaluarse. \\
Reward & Señal que evalúa el éxito y la calidad de la tarea & Determina qué aprenderá el modelo. \\
Replay & Reproducir trayectorias exitosas o fallidas para entrenamiento o análisis & Ayuda a aprender de los errores. \\
Simulation & Entorno controlable o mundo simulado & Reduce el costo de ensayo y error en el mundo real. \\
Memory Update & Escribir la experiencia en la memoria de largo plazo o en un estado externo & Hace que el modelo entienda mejor cuanto más se usa. \\
\bottomrule
\end{tabularx}
\end{knowledgebox}

\subsection{Riesgo: cómo evitar que el aprendizaje en línea caiga en reward hack}

Esta sección añade restricciones. El riesgo del Online Learning es que el modelo aprenda a explotar la recompensa, contamine la memoria, amplifique sesgos, eluda los límites de seguridad o aprenda estrategias erróneas a partir de retroalimentación incorrecta. Por ello, un sistema de aprendizaje en línea realmente utilizable debe contar con permisos por niveles, entornos aislados (sandbox), estados que puedan revertirse, evaluación fuera de línea, auditoría humana y filtros de seguridad.

\begin{warningbox}{El Online Learning no equivale a «dejar que el modelo aprenda lo que sea por su cuenta»}
Si el aprendizaje continuo no tiene límites, provoca contaminación de la memoria, deriva de objetivos y problemas de seguridad. Un Online Learning de alta calidad se parece más a un sistema experimental controlado: tareas claras, recompensas auditables, fallas rastreables y actualizaciones reversibles.
\end{warningbox}

\subsection{Resumen de la sección}

El potencial del Online Learning reside en que convierte el proceso de uso del Agent en un proceso de producción de datos. También trae nuevos problemas de infraestructura y de seguridad. En la siguiente etapa, quien logre operar de forma estable el ciclo cerrado de entornos, recompensas y retroalimentación podría obtener una nueva vía de mejora de los modelos.
